\documentclass[11pt]{article}

\usepackage{amsmath,amssymb,amsthm,mathtools}
\usepackage[margin=1in]{geometry}

\newtheorem{theorem}{Theorem}
\newtheorem{lemma}{Lemma}
\newtheorem{corollary}{Corollary}

\newcommand{\br}[1]{\left(#1\right)}
\newcommand{\cut}{\delta}

\begin{document}

\section*{A valid rounding proof for the one-sided sparsity objective}

Let $G=(V,E,c,w)$ be an instance, where $c\colon E\to\mathbb{R}_{\ge 0}$ is the edge-capacity function and $w\colon V\times V\to\mathbb{R}_{\ge 0}$ is a symmetric demand function. Define
\[
\pi(u)=\sum_{v\in V} w\br{u,v}
\qquad\text{for every }u\in V.
\]
The objective is
\[
 h\br{S}=\frac{c\br{S,\bar S}}{w\br{S,V}}.
\]

\subsection*{1. Augmented instance}

Introduce a dummy vertex $r\notin V$ and set $V^+=V\cup\{r\}$. Define a symmetric demand function $D\colon V^+\times V^+\to\mathbb{R}_{\ge 0}$ by
\[
D\br{u,r}=D\br{r,u}=\pi(u)\qquad\forall u\in V,
\qquad
D\br{u,v}=0\qquad\forall u,v\in V.
\]
Thus every demand pair is of the form $(u,r)$.

For any cut $X\subseteq V^+$, let $S_X$ denote the side of the cut that does not contain $r$. Then
\[
D\br{\cut\br{X}}=\sum_{u\in S_X}\pi(u)=w\br{S_X,V}.
\]
Therefore, minimizing $h(S)$ over $S\subseteq V$ is equivalent to finding a sparsest cut in the augmented instance $(V^+,E,c,D)$, with the convention that the output side is the side not containing $r$.

\subsection*{2. Metric relaxation}

Consider the standard metric relaxation for the augmented sparsest cut instance. Let $d$ be a metric on $V^+$ and normalize so that
\[
\sum_{x,y\in V^+} D\br{x,y}\,d\br{x,y}=1.
\]
The relaxation objective is
\[
\lambda^* = \sum_{uv\in E} c\br{u,v}\,d\br{u,v}.
\]
For an integral cut $X\subseteq V^+$, the associated cut metric is feasible, and the relaxation value is at most the sparsity of the cut. Hence the relaxation is valid.

\subsection*{3. Line-embedding rounding}

We use the standard pairwise-demand line-embedding theorem.

\begin{theorem}[Bourgain-type line embedding]
For every finite metric space $(X,d)$ with $|X|=N$ and every symmetric demand function $D\colon X\times X\to\mathbb{R}_{\ge 0}$, there exists a polynomial-time computable map $f\colon X\to\mathbb{R}$ such that:
\begin{enumerate}
\item $f$ is 1-Lipschitz, i.e.
\[
|f\br{x}-f\br{y}|\le d\br{x,y}\qquad \forall x,y\in X;
\]
\item the average weighted distortion is $O\br{\log N}$, i.e.
\[
\sum_{x,y\in X} D\br{x,y}\,|f\br{x}-f\br{y}|
\ge
\frac{1}{\alpha}
\sum_{x,y\in X} D\br{x,y}\,d\br{x,y}
\]
for some universal constant hidden in $\alpha=O\br{\log N}$.
\end{enumerate}
\end{theorem}

Apply the theorem to the augmented metric space $(V^+,d)$. By translating the line, assume $f\br{r}=0$.

For every threshold $t\in\mathbb{R}$, define the sweep cut
\[
X_t=\{x\in V^+ : f\br{x}\le t\}.
\]

\begin{lemma}[Co-area identity for capacity]
\[
\int_{-\infty}^{\infty} c\br{\cut\br{X_t}}\,dt
=
\sum_{uv\in E} c\br{u,v}\,|f\br{u}-f\br{v}|.
\]
\end{lemma}

\begin{proof}
An edge $uv$ contributes $c\br{u,v}$ to the cut $\cut\br{X_t}$ precisely for those thresholds $t$ lying between $f\br{u}$ and $f\br{v}$. Thus the total contribution of that edge to the integral is $c\br{u,v}\,|f\br{u}-f\br{v}|$. Summing over all edges gives the identity.
\end{proof}

\begin{lemma}[Co-area identity for demand]
\[
\int_{-\infty}^{\infty} D\br{\cut\br{X_t}}\,dt
=
\sum_{x,y\in V^+} D\br{x,y}\,|f\br{x}-f\br{y}|.
\]
\end{lemma}

\begin{proof}
The proof is identical. For each demand pair $xy$, the pair contributes $D\br{x,y}$ to the cut whenever $t$ lies between $f\br{x}$ and $f\br{y}$, and therefore contributes exactly $D\br{x,y}\,|f\br{x}-f\br{y}|$ to the integral.
\end{proof}

Since $f$ is 1-Lipschitz,
\[
\int_{-\infty}^{\infty} c\br{\cut\br{X_t}}\,dt
= 
\sum_{uv\in E} c\br{u,v}\,|f\br{u}-f\br{v}|
\le
\sum_{uv\in E} c\br{u,v}\,d\br{u,v}
= \lambda^*.
\]
By the embedding theorem,
\[
\int_{-\infty}^{\infty} D\br{\cut\br{X_t}}\,dt
=
\sum_{x,y\in V^+} D\br{x,y}\,|f\br{x}-f\br{y}|
\ge
\frac{1}{\alpha}
\sum_{x,y\in V^+} D\br{x,y}\,d\br{x,y}
=
\frac{1}{\alpha},
\]
where the last equality is the normalization of the relaxation.

Suppose, for contradiction, that every threshold $t$ satisfies
\[
\frac{c\br{\cut\br{X_t}}}{D\br{\cut\br{X_t}}} > \alpha\lambda^*.
\]
Then integrating over all thresholds gives
\[
\int_{-\infty}^{\infty} c\br{\cut\br{X_t}}\,dt
>
\alpha\lambda^*
\int_{-\infty}^{\infty} D\br{\cut\br{X_t}}\,dt
\ge
\alpha\lambda^*\cdot \frac{1}{\alpha}
=
\lambda^*,
\]
contradicting the upper bound on the capacity integral. Therefore there exists a threshold $t^*$ such that
\[
\frac{c\br{\cut\br{X_{t^*}}}}{D\br{\cut\br{X_{t^*}}}}
\le
\alpha\lambda^*.
\]
Since the cut changes only when $t$ crosses one of the finitely many values $f\br{x}$, it suffices to check the $|V^+|-1$ prefix cuts. This is polynomial time.

\subsection*{4. Returning to the original objective}

Let $S$ be the side of the chosen cut $X_{t^*}$ that does not contain $r$. By construction,
\[
D\br{\cut\br{X_{t^*}}}=w\br{S,V}
\qquad\text{and}\qquad
c\br{\cut\br{X_{t^*}}}=c\br{S,\bar S}.
\]
Hence
\[
 h\br{S}=
\frac{c\br{S,\bar S}}{w\br{S,V}}
=
\frac{c\br{\cut\br{X_{t^*}}}}{D\br{\cut\br{X_{t^*}}}}
\le
\alpha\lambda^*.
\]
Because $\alpha=O\br{\log n}$, this gives an $O\br{\log n}$-approximation algorithm for the original objective.

\begin{corollary}
The proof is fully formal: every lower bound comes from the pairwise-demand embedding theorem applied to the augmented instance, and no unsupported one-sided embedding statement is used.
\end{corollary}

\end{document}
